\documentclass[11pt,a4paper]{article}
\usepackage[margin=1in]{geometry}
\usepackage{amsmath,amssymb}
\title{Topology-Aware Adaptive Optimization for Graph Learning}
\author{Open-Source Research Note}
\date{}

\begin{document}
\maketitle

\begin{abstract}
We summarize a topology-aware variant of Adam where effective learning rates depend on graph structural indicators. The objective is to improve optimization stability across heterogeneous graph regimes.
\end{abstract}

\section{Problem Statement}
For parameters $\theta$ trained on graph-structured data $G=(V,E)$, standard Adam updates are
\[
\theta_{t+1}=\theta_t-\eta\,\frac{\hat{m}_t}{\sqrt{\hat{v}_t}+\epsilon}.
\]
We seek graph-conditioned scaling $s_t(G)$ and optional node-local weights $w_v(G)$ to modulate updates:
\[
\theta_{t+1}=\theta_t-\eta\,s_t(G)\,\frac{\hat{m}_t}{\sqrt{\hat{v}_t}+\epsilon}.
\]

\section{Formalization}
Define a topological relevance function (TRF),
\[
\mathrm{TRF}(G)=1+\mu_G\,\Phi(G),
\]
where $\Phi$ combines normalized structural statistics (e.g., connectivity, efficiency, path length). Global scaling may use
\[
s_t(G)=\mathrm{TRF}(G)^{-1}.
\]
For local scaling on node-indexed gradients, define $w_v\propto g(\deg(v),\mathrm{homophily}(v),\ldots)$ and apply weighted aggregation before the optimizer step.

\section{Research Contribution}
The project offers an open-source implementation and synthetic graph benchmarks to study when topological adaptation modifies convergence speed and terminal loss versus standard adaptive optimizers.

\begin{thebibliography}{9}
\bibitem{kingma2015} D. P. Kingma, J. Ba, ``Adam: A Method for Stochastic Optimization,'' \emph{ICLR}, 2015.
\bibitem{reddi2018} S. Reddi et al., ``On the Convergence of Adam and Beyond,'' \emph{ICLR}, 2018.
\bibitem{kipf2017} T. Kipf, M. Welling, ``Semi-Supervised Classification with Graph Convolutional Networks,'' \emph{ICLR}, 2017.
\bibitem{newman2010} M. Newman, \emph{Networks: An Introduction}, Oxford Univ. Press, 2010.
\bibitem{paszke2019} A. Paszke et al., ``PyTorch: An Imperative Style, High-Performance Deep Learning Library,'' \emph{NeurIPS}, 2019.
\end{thebibliography}

\end{document}
